\section{Reset-free 控制与记忆管理}
\subsection{重置假设的弱化}
传统 RL 依赖人为或 simulator 重置。Chelsea 提出“让 policy 本身来 reset”——在每轮任务结束后启动另一个 policy，让环境回到新的初始集合，使 robot 能持续运行。

\begin{figure}[H]
\centering
\includegraphics[width=0.88\textwidth]{slides-images/slide-04.jpg}
\caption{Slide 4 讲解前向/后向控制器如何交替执行任务与 reset。}
\end{figure}

\noindent 图中 timeline 表示一次任务执行后的回收过程：任务 policy（中间）完成 goal 后，将 current state 交给 reset policy。reset policy 运行后把 system 带回 base manifold，继续提供任务起点。

\begin{importantbox}{Forward-Backward Controller 机制}
\begin{itemize}
\item $\pi_f$: 从 initial state 探索 task-specific trajectory；\\
\item $\pi_b$: 接受当前 state，尝试回到 initial manifold；\\
\item 轮流训练两者，$\pi_b$ 的成功率定义新 episode 的起点。
\end{itemize}
\end{importantbox}

\subsection{持续 episodic buffer}
Reset-free 要求 long-lived buffer 记住“最近 failed states”，以便策略不断学习重置。Chelsea 建议维护 prioritized replay buffer，同时引入\textbf{time-aware sampling}，优先回放发生在最近一次 reset 之后的数据。

\begin{knowledgebox}{Time-Aware Replay}
\begin{itemize}
\item 将每个 transition 附加 timestamp；\\
\item 在训练中采用 $w_t \propto \exp(-\lambda (t_{now}-t))$ 近似 latest priority；\\
\item 保证新数据被充分利用，同时保留少量 rare event。
\end{itemize}
\end{knowledgebox}

\subsection{经验撤销与环境多样性}
讲者进一步指出：自主系统需具备\textbf{撤销经验}的能力，即模型在 fail 的轨迹上回退，并尝试新的策略组合。这涉及动态规划多个 reset-policy 以覆盖不同物理扰动。
\begin{warningbox}{Fail-fast but recover-fast}
鼓励策略快速尝试多个动作组合，如果当前控制失败，马上启动 reset policy；同时记录失败模式供 future policy adaptation。
\end{warningbox}

\subsection{本章小结}
Reset-free 框架用 forward-backward controllers、time-aware buffer 和 fast recovery 实现持续运行，避免依赖人工 reset。
\newpage

